% =========================================================================
\section{Simulaciones: tetrámero de cuatro esferas}
\label{sec:tetramero}
% =========================================================================

Se extiende el análisis al caso de cuatro esferas en configuración lineal,
con el eje de la cadena orientado a lo largo de $x$ y la onda incidente
propagándose a lo largo de $z$. Se consideran cuatro configuraciones:
el caso homogéneo con $a=25\,\text{nm}$, el caso homogéneo con
$a=85\,\text{nm}$, y dos configuraciones heterogéneas que combinan ambos
tamaños con distinta disposición geométrica. Como en el estudio del dímero 
el gap más estrecho produce la mayor intensificación, las simulaciones se 
realizan únicamente para $d=5\,\text{nm}$.

% =========================================================================
\subsection{Tetrámero homogéneo: $a = 25\,\text{nm}$}
\label{subsec:tetramero_a25}
% =========================================================================
Este caso sirve de para analizar cómo cambia el acoplamiento
al añadir más partículas respecto al dímero equivalente.

\begin{figure}[H]
  \centering
  \begin{subfigure}[b]{0.48\textwidth}
    \centering
    \includegraphics[width=\textwidth]{../../Simulaciones/Definitivo/simulacion-de-referencia/4e/ew/a25/sep5/sw/plots/Eficiencias_-_Unpol.png}
    \caption{Eficiencias espectrales. Pico en $\lambda\approx550\,\text{nm}$
             con $Q_\text{abs}$ dominante.}
    \label{fig:tet_a25_eff}
  \end{subfigure}
  \hfill
  \begin{subfigure}[b]{0.48\textwidth}
    \centering
    \includegraphics[width=\textwidth]{../../Simulaciones/Definitivo/simulacion-de-referencia/4e/ew/a25/sep5/wl550/plots/Figure_1.png}
    \caption{Campo cercano $|\mathbf{E}_\text{avg}|^2$, $\lambda=550\,\text{nm}$.
             Máximo: $430.78\,|\mathbf{E}_0|^2$.}
    \label{fig:tet_a25_nf}
  \end{subfigure}
  \caption{Tetrámero homogéneo con $a=25\,\text{nm}$ y $d=5\,\text{nm}$.
           La resonancia se mantiene en torno a $\lambda\approx550\,\text{nm}$
           con absorción dominante. La zona de concentración se localiza
           en el gap central, entre las dos esferas interiores, mientras
           que los gaps exteriores muestran una intensificación
           significativamente menor.}
  \label{fig:tet_a25}
\end{figure}

El tetrámero con $a=25\,\text{nm}$ opera en el mismo régimen que el dímero
equivalente: la resonancia no se desplaza y la absorción sigue siendo el
canal dominante. El efecto más notable es la concentración del campo en el
gap central, que alcanza $430.78\,|\mathbf{E}_0|^2$, notablemente superior
a los $147\,|\mathbf{E}_0|^2$ del dímero, lo que indica que añadir más
partículas intensifica el campo en los gaps interiores de la cadena.


% =========================================================================
\subsection{Tetrámero homogéneo: $a = 85\,\text{nm}$}
\label{subsec:tetramero_a85}
% =========================================================================

\begin{figure}[H]
  \centering
  \begin{subfigure}[b]{0.32\textwidth}
    \centering
    \includegraphics[width=\textwidth]{../../Simulaciones/Definitivo/simulacion-de-referencia/4e/ew/a85/sep5/sw/plots/Eficiencias_-_Unpol.png}
    \caption{Eficiencias espectrales.}
    \label{fig:tet_a85_eff}
  \end{subfigure}
  \hfill
  \begin{subfigure}[b]{0.32\textwidth}
    \centering
    \includegraphics[width=\textwidth]{../../Simulaciones/Definitivo/simulacion-de-referencia/4e/ew/a85/sep5/wl590/plots/Figure_1.png}
    \caption{Campo cercano, $\lambda=590\,\text{nm}$.
             Máximo: $131.41\,|\mathbf{E}_0|^2$.}
    \label{fig:tet_a85_nf_590}
  \end{subfigure}
  \hfill
  \begin{subfigure}[b]{0.32\textwidth}
    \centering
    \includegraphics[width=\textwidth]{../../Simulaciones/Definitivo/simulacion-de-referencia/4e/ew/a85/sep5/wl1000/plots/Figure_1.png}
    \caption{Campo cercano, $\lambda=1000\,\text{nm}$.
             Máximo: $323.74\,|\mathbf{E}_0|^2$.}
    \label{fig:tet_a85_nf_1000}
  \end{subfigure}
  \caption{Tetrámero homogéneo con $a=85\,\text{nm}$ y $d=5\,\text{nm}$.
           El espectro muestra tres picos: uno dominado por la absorción
           en torno a $\lambda\approx500\,\text{nm}$ y dos por la dispersión
           en torno a $\lambda\approx590\,\text{nm}$ y $\lambda\approx1000\,\text{nm}$.
           A $\lambda=590\,\text{nm}$ la zona de concentración
           ($131.41\,|\mathbf{E}_0|^2$) se localiza en los gaps exteriores.
           A $\lambda=1000\,\text{nm}$ el máximo ($323.74\,|\mathbf{E}_0|^2$)
           se desplaza al gap central, con los gaps exteriores mostrando
           una intensificación menor.}
  \label{fig:tet_a85}
\end{figure}

El acoplamiento entre las cuatro partículas desplaza el pico principal
hasta el infrarrojo cercano, muy por encima de lo observado en el dímero
equivalente. El campo cercano revela además que los gaps activos cambian
con la longitud de onda: a $590\,\text{nm}$ dominan los gaps exteriores,
mientras que a $1000\,\text{nm}$ es el gap central el que concentra
la mayor intensidad.

% =========================================================================
\subsection{Tetrámero heterogéneo alternado}
\label{subsec:tetramero_alterno}
% =========================================================================

\begin{figure}[H]
  \centering
  \begin{subfigure}[b]{0.48\textwidth}
    \centering
    \includegraphics[width=\textwidth]{../../Simulaciones/Definitivo/simulacion-de-referencia/4e/ew/a85a25/alterno/sw/plots/Eficiencias_-_Unpol.png}
    \caption{Eficiencias espectrales.}
    \label{fig:tet_alt_eff}
  \end{subfigure}
  \hfill
  \begin{subfigure}[b]{0.48\textwidth}
    \centering
    \includegraphics[width=\textwidth]{../../Simulaciones/Definitivo/simulacion-de-referencia/4e/ew/a85a25/alterno/wl660/plots/Figure_1.png}
    \caption{Campo cercano, $\lambda=660\,\text{nm}$.
             Máximo: $305.53\,|\mathbf{E}_0|^2$.}
    \label{fig:tet_alt_nf}
  \end{subfigure}
  \caption{Tetrámero heterogéneo alternado
            ($85$--$25$--$85$--$25\,\text{nm}$) con $d=5\,\text{nm}$.
            El espectro muestra dos picos: uno en torno a
            $\lambda\approx520\,\text{nm}$ donde $Q_\text{abs}$ todavía
            contribuye, y un pico dispersor principal en
            $\lambda\approx660\,\text{nm}$. La zona de concentración se
            localiza en torno a la esfera pequeña flanqueada por las dos
            grandes, con un máximo de $305.53\,|\mathbf{E}_0|^2$;
            el gap exterior permanece significativamente menos activo.}
  \label{fig:tet_alt}
\end{figure}

FALTA ALGO

% =========================================================================
\subsection{Tetrámero heterogéneo con esferas pequeñas centrales}
\label{subsec:tetramero_central}
% =========================================================================

\begin{figure}[H]
  \centering
  \begin{subfigure}[b]{0.32\textwidth}
    \centering
    \includegraphics[width=\textwidth]{../../Simulaciones/Definitivo/simulacion-de-referencia/4e/ew/a85a25/medio/sw/plots/Eficiencias_-_Unpol.png}
    \caption{Eficiencias espectrales.}
    \label{fig:tet_medio_eff}
  \end{subfigure}
  \hfill
  \begin{subfigure}[b]{0.32\textwidth}
    \centering
    \includegraphics[width=\textwidth]{../../Simulaciones/Definitivo/simulacion-de-referencia/4e/ew/a85a25/medio/wl520/plots/Figure_1.png}
    \caption{Campo cercano, $\lambda=520\,\text{nm}$.
             Máximo: $109.63\,|\mathbf{E}_0|^2$.}
    \label{fig:tet_medio_nf_520}
  \end{subfigure}
  \hfill
  \begin{subfigure}[b]{0.32\textwidth}
    \centering
    \includegraphics[width=\textwidth]{../../Simulaciones/Definitivo/simulacion-de-referencia/4e/ew/a85a25/medio/wl620/plots/Figure_1.png}
    \caption{Campo cercano, $\lambda=620\,\text{nm}$.
             Máximo: $566.20\,|\mathbf{E}_0|^2$.}
    \label{fig:tet_medio_nf_620}
  \end{subfigure}
  \caption{Tetrámero heterogéneo ($85$--$25$--$25$--$85\,\text{nm}$) con
           $d=5\,\text{nm}$. El espectro muestra un pico de extinción en
           $\lambda\approx520\,\text{nm}$ dominado por la absorción, y un
           pico dispersor principal en $\lambda\approx620\,\text{nm}$.
           En ambos casos la zona de concentración se localiza entre las
           dos esferas pequeñas centrales; a $\lambda=620\,\text{nm}$
           la intensificación alcanza $566.20\,|\mathbf{E}_0|^2$, el valor
           más alto de todas las configuraciones estudiadas.}
  \label{fig:tet_medio}
\end{figure}

Situar las dos esferas pequeñas en el interior concentra el gap más estrecho
en la zona central de la cadena, donde el acoplamiento es máximo. El resultado
es la mayor intensificación de campo observada en este trabajo: $566.20\,|\mathbf{E}_0|^2$
a $\lambda=620\,\text{nm}$, frente a los $430.78\,|\mathbf{E}_0|^2$ del
tetrámero homogéneo de $25\,\text{nm}$ y los $305.53\,|\mathbf{E}_0|^2$
de la configuración alternada.